% Generated by tex/build_swebok_structure.py from tex/chapters/ch01/03_ソフトウェア要求の基礎.tex
\subsubsection{派生要求}

\term{派生要求}とは、外部の利害関係者から直接出された要求ではなく、開発チーム内部の判断から生まれ、下位チームや構成要素に課される要求である。

たとえば、アーキテクトが「パイプ・アンド・フィルタ構造を採用する」と決めたとする。この決定は、外部顧客から見ると設計判断である。しかし、個々のフィルタを作るサブチームから見ると、「この形式で入力を受け取り、この形式で出力する」という要求になる。このように、要求か設計かは視点に依存する。

\MiniBox{視点依存性}{同じ記述でも、上位の視点では設計判断、下位の視点では要求になることがある。要求とは「その作業を始める人にとって、満たさなければならない入力条件」として理解するとよい。}
